\section{Conclusion}\label{sec:conclusion}

For \((1,s)\)-nucleus decomposition, CPI paths do not need to be mutable residual state.
%
Vertex peeling removes vertices but does not delete edges, so the path structure built at the beginning remains valid throughout the peel.
%
A hold removal kills a path contribution.
%
A pivot removal decreases an effective pivot counter.
%
This changes exact decomposition from residual-index maintenance to counter maintenance over a static symbolic index.

The algorithmic consequences follow from that state change.
%
\localH{} states the fixed-point computation over the static \cpi{}.
%
\peelH{} realizes the same computation in deletion order.
%
\paraBuild{} addresses construction after cheap peeling exposes it as the next bottleneck.
%
\buildHier{} reads the sorted deletion log in reverse to recover the hierarchy.

The experiments and case studies support the practical value of the specialization.
%
The reported matched configurations agree with the mutable baseline, and the measured pipeline reduces runtime and memory on the evaluated benchmarks.
%
The case studies show that vertex-centric \((1,s)\) can match higher-\(r\) quality metrics on the reported community tasks at much lower cost.
%
The boundary is explicit: the static-\cpi{} invariant is for \(r=1\), not for general \((r,s)\).

A natural next step is to identify restricted \(r\ge2\) settings where enough path structure survives to support counter-only updates.
